\begin{center}
{\large\bfseries\color{azulPDMI}Quanto precisa de virar o gato?}\\[2pt]
{\small Resolução das questões 1.1 a 4.4, conclusão e desafio facultativo}
\end{center}

\begin{nota}
\textbf{Condições consideradas.} Estuda-se o contorno de polígonos
\textbf{convexos}, percorrido uma única vez, sem cruzamentos e sempre
no mesmo sentido. Em cada vértice considera-se \textbf{um} ângulo
externo, correspondente à mudança de direção do gato. O programa
faz uma viragem depois de cada lado, \textbf{incluindo o último}.
Mover o gato com a caneta levantada para o centro do palco não integra
o percurso geométrico.
\end{nota}

\section*{Parte I}

\noindent\textbf{1.1. Quatro lados e uma viragem de $80^\circ$}

Não se obtém um quadrado. Num quadrado, os lados consecutivos são
perpendiculares, pelo que a mudança de direção em cada vértice tem
de ser de $90^\circ$. Se o gato virar apenas $80^\circ$, a direção
do lado seguinte não será perpendicular à do lado anterior. A
sequência de quatro segmentos de igual comprimento não fecha o
contorno de um quadrado.

A soma das viragens programadas é
\[
4\times80^\circ=320^\circ,
\]
isto é, $40^\circ$ menos do que uma volta completa. Esta conta ajuda
a interpretar o resultado, mas a razão por que o contorno de um
polígono convexo corresponde a uma volta completa será justificada
na questão~4.4; não deve ser simplesmente pressuposta no lançamento
da tarefa.

\begin{center}
\begin{minipage}{.46\linewidth}
\centering
\begin{tikzpicture}[scale=.72,>=Stealth]
  \draw[thick,azulPDMI,line join=round]
  (0,0)--(2,0)--(2.347,-1.970)--(.468,-2.654)--(-.532,-.922);
  \fill[azulPDMI] (0,0) circle (2.5pt);
  \fill[orange!85!black] (-.532,-.922) circle (2.5pt);
  \node[above left] at (0,0) {$A$};
  \node[left] at (-.532,-.922) {$A'$};
  \draw[gray,dashed] (-.532,-.922)--(0,0);
\end{tikzpicture}\\[-3pt]
{\footnotesize $4$ viragens de $80^\circ$: $A'\neq A$.}
\end{minipage}\hfill
\begin{minipage}{.46\linewidth}
\centering
\begin{tikzpicture}[scale=.72,>=Stealth]
  \draw[thick,azulPDMI,line join=round]
    (0,0)--(2,0)--(2,-2)--(0,-2)--cycle;
  \fill[azulPDMI] (0,0) circle (2.5pt);
  \node[above left] at (0,0) {$A$};
\end{tikzpicture}\\[-3pt]
{\footnotesize $4$ viragens de $90^\circ$: quadrado.}
\end{minipage}
\end{center}

\noindent\textbf{1.2. Viragem necessária para desenhar o quadrado}

A amplitude correta é $\boxed{90^\circ}$. Como o gato anda sempre
o mesmo número de passos, os quatro lados têm o mesmo comprimento.
Ao realizar uma viragem de $90^\circ$ depois de cada lado, obtêm-se
quatro ângulos retos e o contorno fecha sem cruzamentos. A soma das
quatro mudanças de direção é
\[
90^\circ+90^\circ+90^\circ+90^\circ=360^\circ.
\]

\noindent\textbf{1.3. Comparação entre ângulo interno e viragem}

A coincidência verificada no quadrado \textbf{não é geral}. Num
triângulo equilátero, por exemplo, cada ângulo interno mede $60^\circ$,
mas o gato tem de virar $120^\circ$ ao passar de um lado ao seguinte.
Os dois ângulos são adjacentes e suplementares:
\[
60^\circ+120^\circ=180^\circ.
\]
Num pentágono regular, o ângulo interno mede $108^\circ$ e a
viragem mede $72^\circ$. A igualdade entre a viragem e o ângulo
interno no quadrado é, portanto, um caso particular.

\begin{nota}
\textbf{Resposta possível de um aluno.} «Não. No quadrado os dois
ângulos medem $90^\circ$, mas num triângulo equilátero o gato vira
$120^\circ$ e o ângulo interno mede $60^\circ$. Juntos perfazem
$180^\circ$.»
\end{nota}

\section*{Parte II}

\noindent\textbf{2.1. Tabela completa}

\begin{center}
\renewcommand{\arraystretch}{1.32}
\begin{tabularx}{\linewidth}{|>{\raggedright\arraybackslash}X|
>{\centering\arraybackslash}p{1.5cm}|
>{\centering\arraybackslash}p{2.6cm}|
>{\centering\arraybackslash}p{2.65cm}|}
\hline
\rowcolor{azulClaro}
\textbf{Polígono} & \textbf{N.º de lados} &
\textbf{Viragem por vértice} & \textbf{Soma das viragens}\\\hline
Triângulo equilátero & $3$ & $120^\circ$ & $3\times120^\circ=360^\circ$\\\hline
Quadrado & $4$ & $90^\circ$ & $4\times90^\circ=360^\circ$\\\hline
Pentágono regular & $5$ & $72^\circ$ & $5\times72^\circ=360^\circ$\\\hline
Hexágono regular & $6$ & $60^\circ$ & $6\times60^\circ=360^\circ$\\\hline
\end{tabularx}
\end{center}

Em cada caso, o valor indicado é a \textbf{mudança de direção},
não a amplitude de um ângulo interno. Os valores correspondem ao
percurso do contorno do polígono regular convexo, sem cruzamentos e
com uma viragem em cada vértice. A mera igualdade
$n\,v=360^\circ$ não garante, sem outras condições, que um
percurso arbitrário seja o contorno simples pretendido.

\noindent\textbf{2.2. O que muda e o que se mantém?}

À medida que o número de lados aumenta de $3$ para $6$, aumenta o
número de viragens, mas diminui a amplitude de cada uma: $120^\circ$,
$90^\circ$, $72^\circ$ e $60^\circ$. A \textbf{soma} das amplitudes das
viragens mantém-se igual a $360^\circ$. Nos polígonos regulares,
todas as viragens de um mesmo percurso têm a mesma amplitude.

É importante distinguir três quantidades: \emph{número de lados},
\emph{amplitude de cada viragem} e \emph{soma das viragens}. Esta
última é a quantidade que se mantém nos casos estudados.

\noindent\textbf{2.3. Generalização para $n$ lados}

Se o gato percorre uma única vez o contorno de um polígono regular
convexo com $n$ lados, efetuando uma viragem de amplitude $v$ em cada
vértice, efetua $n$ viragens iguais. Como perfaz uma volta completa,
\[
\underbrace{v+v+\cdots+v}_{n\text{ parcelas}}=360^\circ
\quad\Longleftrightarrow\quad
n\,v=360^\circ.
\]
Logo,
\[
\boxed{v=\frac{360^\circ}{n}},\qquad n\in\mathbb{N},\ n\geq3.
\]
Esta relação dá a amplitude de \textbf{cada} ângulo externo de um
\textbf{polígono regular convexo}. Não significa que todos os ângulos
externos de um polígono convexo qualquer sejam iguais.

\section*{Parte III}

\noindent\textbf{3.1. Viragens no paralelogramo}

Ao premir a tecla \textbf{I}, o gato desenha um paralelogramo não
retângulo com quatro mudanças de direção, por esta ordem:
\[
60^\circ,\quad120^\circ,\quad60^\circ,\quad120^\circ.
\]
A sua soma é
\[
60^\circ+120^\circ+60^\circ+120^\circ=\boxed{360^\circ}.
\]
As amplitudes dos ângulos internos correspondentes são,
respetivamente, $120^\circ$, $60^\circ$, $120^\circ$ e $60^\circ$.

\noindent\textbf{3.2. Comparação e alcance da experiência}

Sim, a soma é igual à obtida para os polígonos regulares: $360^\circ$.
O paralelogramo mostra que \textbf{a igualdade das amplitudes das
viragens não é necessária} para se obter essa soma: há duas viragens
de $60^\circ$ e duas de $120^\circ$.

Contudo, \textbf{um paralelogramo não demonstra} que a propriedade
é válida para todos os polígonos convexos. A experiência fornece
um novo exemplo favorável à conjetura, mas é necessário apresentar
um argumento matemático que abranja qualquer polígono convexo.

\begin{nota}
\textbf{Resposta possível de um aluno.} «A soma continua a ser
$360^\circ$, apesar de o gato não virar sempre o mesmo. Mas testar
este paralelogramo não chega para saber o que acontece com todos os
polígonos convexos.»
\end{nota}

\section*{Parte IV}

\noindent\textbf{4.1. Relação entre o ângulo interno e o externo}

O ângulo externo $\beta$ considerado na tarefa é formado pelo
prolongamento de um lado do polígono e pelo lado seguinte. O ângulo
interno $\alpha$ e esse ângulo externo têm um lado comum; os
outros lados são semirretas opostas. São, portanto, \textbf{ângulos
adjacentes suplementares}, pelo que
\[
\boxed{\alpha+\beta=180^\circ}
\qquad\text{ou}\qquad
\boxed{\alpha=180^\circ-\beta}.
\]
A mudança de direção do gato corresponde ao ângulo externo,
\textbf{não}, em geral, ao ângulo interno do polígono.

\noindent\textbf{4.2. O pentágono convexo}

Um pentágono tem cinco vértices e, por isso, cinco pares de um
ângulo interno com o respetivo externo adjacente. Cada par soma
$180^\circ$; assim, a soma das amplitudes dos \textbf{dez ângulos}
(cinco internos e cinco externos) é
\[
5\times180^\circ=900^\circ.
\]
A soma dos cinco ângulos externos é $360^\circ$. Logo, a soma dos
cinco ângulos internos é
\[
\boxed{S_5=900^\circ-360^\circ=540^\circ}.
\]
A resposta não depende de o pentágono ser regular. Se fosse regular,
teríamos cinco ângulos internos iguais a $108^\circ$ e verificar-se-ia
$5\times108^\circ=540^\circ$.

\noindent\textbf{4.3. Generalização para um polígono com $n$ lados}

Considere-se um polígono convexo com $n\geq3$ lados. Designemos por
$\alpha_i$ a amplitude do ângulo interno do vértice $i$ e por
$\beta_i$ a amplitude do ângulo externo correspondente, escolhido
sempre no mesmo sentido do percurso. Como cada par é suplementar,
\[
\alpha_i+\beta_i=180^\circ
\qquad(i=1,\ldots,n).
\]
Somando estas $n$ igualdades, obtém-se
\[
\underbrace{\sum_{i=1}^{n}\alpha_i}_{S}
+\underbrace{\sum_{i=1}^{n}\beta_i}_{360^\circ}
=n\times180^\circ.
\]
Consequentemente,
\[
\boxed{S=n\times180^\circ-360^\circ=(n-2)\times180^\circ}.
\]
Na expressão, $n\times180^\circ$ é a soma das amplitudes dos
$n$ pares de ângulos suplementares; subtraem-se $360^\circ$ porque
essa é a soma das amplitudes dos ângulos externos. O resultado é a
soma das amplitudes dos ângulos \textbf{internos} do polígono.

Esta dedução não exige que os ângulos internos sejam todos iguais;
por isso, aplica-se aos polígonos convexos regulares e não regulares.

\noindent\textbf{4.4. Justificação geométrica da soma dos ângulos externos}

A explicação não deve limitar-se à afirmação de que o gato
«regressou ao sítio de onde partiu». Na verdade, a orientação final
poderia coincidir com a inicial após zero, uma ou várias voltas
completas; e um desenho poderia fechar com autointerseções.

A propriedade resulta das condições geométricas do percurso:
\begin{enumerate}[label=\textbf{(\roman*)},leftmargin=1.15cm,itemsep=2pt,topsep=3pt]
\item o gato percorre \textbf{uma única vez} o contorno de um polígono
\textbf{convexo}, sem o atravessar;
\item em cada vértice muda de direção \textbf{sempre no mesmo sentido},
pois todos os ângulos externos considerados são viragens exteriores;
\item ao acompanhar, por ordem, as orientações dos lados de um
contorno convexo, a direção percorre as orientações de uma
\textbf{única volta completa}, sem retroceder nem fazer uma segunda volta.
\end{enumerate}

Assim, a soma das mudanças de direção, incluindo a última viragem
que recupera a orientação inicial, é
\[
\boxed{\beta_1+\beta_2+\cdots+\beta_n=360^\circ}.
\]

\textbf{Uma forma de tornar o argumento visível na discussão:}
imaginar os vetores orientados de cada lado transportados, sem
rodar, para um mesmo ponto. Como o contorno é convexo, as direções
sucedem-se pela mesma ordem e os ângulos entre direções consecutivas
preenchem exatamente uma volta em torno desse ponto. Não há
sobreposições de voltas nem direções percorridas em sentido inverso.
Este argumento geométrico explica a regularidade para \textbf{qualquer}
polígono convexo, em vez de depender de uma lista finita de testes.

\begin{nota}
\textbf{Precisão matemática.} A afirmação de que as direções dos
lados perfazem exatamente uma volta usa a \emph{convexidade} e o
facto de o contorno ser percorrido uma vez. Não decorre apenas de o
gato terminar na mesma posição ou orientação. Uma justificação
mais formal pode recorrer à variação total de direção num contorno
convexo; ao nível do 7.º ano, o argumento por uma volta completa,
apoiado na representação das direções, é a explicação geométrica
adequada a explorar e discutir.
\end{nota}

\section*{Conclusão a registar na ficha}

\begin{conclusao}
Para um polígono convexo com $n$ lados, a soma das amplitudes dos
seus ângulos internos é
\[
\boxed{S=(n-2)\times180^\circ}.
\]
\textbf{Justificação:} em cada vértice, o ângulo interno e o externo
adjacente somam $180^\circ$; existem $n$ pares desse tipo. Ao
percorrer uma vez o contorno de um polígono convexo, a soma dos
ângulos externos é $360^\circ$. Por isso, a soma dos ângulos
internos obtém-se subtraindo $360^\circ$ a $n\times180^\circ$.
\end{conclusao}

\Needspace{14\baselineskip}
\section*{Desafio facultativo}

O gato desenha uma \textbf{estrela de cinco pontas} (pentagrama),
em que os segmentos se cruzam. Ao fim dos cinco segmentos, volta
à posição inicial, mas o desenho \textbf{não é o contorno simples
de um pentágono convexo}. A soma das cinco viragens programadas é
\[
5\times144^\circ=\boxed{720^\circ}=2\times360^\circ.
\]
O gato completou \textbf{duas voltas na orientação}, embora tenha
regressado ao ponto de partida. Logo, fechar um percurso não é
condição suficiente para garantir que a soma das viragens seja
$360^\circ$. É indispensável explicitar as condições de
\emph{convexidade}, \emph{contorno simples} e \emph{uma única volta}.

\begin{center}
\begin{tikzpicture}[scale=.87,line join=round]
\foreach \k in {0,...,4}{
  \coordinate (V\k) at ({90-72*\k}:1.85);
}
\draw[very thick,azulPDMI]
 (V0)--(V2)--(V4)--(V1)--(V3)--cycle;
\fill[azulPDMI] (V0) circle (2.2pt);
\node[above] at (V0) {Início / fim};
\end{tikzpicture}
\end{center}

\section*{Síntese da resolução e cuidados na discussão}

\begin{enumerate}[label=\textbf{\arabic*.},leftmargin=1.05cm,itemsep=3pt,topsep=3pt]
\item \textbf{O que o Scratch permite observar:} valores de viragem
para polígonos regulares e um exemplo não regular em que a soma
exterior continua a ser $360^\circ$.
\item \textbf{O que precisa de ser justificado:} por que razão a
mudança total de direção, num único percurso do contorno convexo,
é uma volta completa.
\item \textbf{A passagem essencial:} a relação
$\alpha_i+\beta_i=180^\circ$ em cada um dos $n$ vértices, seguida
da subtração da soma dos ângulos externos.
\item \textbf{Um erro a prevenir:} escrever $v=360^\circ/n$
para \emph{cada} ângulo externo de um polígono não regular. Essa
igualdade só é válida quando as viragens são todas iguais.
\item \textbf{Outro erro a prevenir:} concluir que a soma é $360^\circ$
por o gato ter regressado ao ponto inicial; a estrela mostra por
que esse argumento é insuficiente.
\item \textbf{Articulação com a primeira aula:} chega-se à mesma
expressão $S=(n-2)\times180^\circ$, mas por ângulos externos e
suplementaridade, \textbf{sem decompor o polígono em triângulos}.
\end{enumerate}